\section{Conclusion}
In this work, we investigate the coupling between weight decay and other optimization hyperparameters from a generalization perspective. By establishing generalization bounds, we show that introducing the weight decay coefficient $\lambda$ disrupts the pre-existing hyperparameter equilibrium, necessitating a system-wide readjustment. Based on this, we propose the "Regularization Equivalence" hypothesis as a heuristic design principle. By aligning the stability profiles of explicit regularization (weight decay) and implicit regularization (SWA), we derive the explicit hyperparameter scaling law: $\lambda \propto 1/\eta T$. 
Furthermore, our analysis explicitly characterizes the role of batch size $B$ as a noise compensation mechanism: as $B$ increases, the reduction in implicit gradient noise must be counterbalanced by scaling up the joint configuration of $\lambda$ and $\eta$.
Empirically, we validated our heuristic theoretical findings and the proposed scaling laws across diverse architectures, ranging from ResNet-18 to Qwen3-0.6B language model. 



%Our current theoretical analysis is grounded in convex assumptions. In future work, we aim to extend this framework to non-convex settings, thereby providing solid theoretical guidance for hyperparameter configuration in large-scale deep learning training.